% !TEX program = pdflatex
\documentclass[11pt]{article}
\usepackage[margin=0.8in]{geometry}
\usepackage{amsmath,amssymb}
\usepackage{enumitem}
\usepackage{titlesec}
\usepackage{xcolor}

\titleformat{\section}{\large\bfseries}{}{0em}{}[\vspace{-0.3em}\rule{\textwidth}{0.4pt}]
\setlength{\parindent}{0pt}
\setlength{\parskip}{0.4em}

\definecolor{qcolor}{HTML}{1565C0}
\newcommand{\Q}[1]{\textcolor{qcolor}{\textbf{Q: #1}}}

\begin{document}

{\Large\textbf{Week 4 Meeting Cheat Sheet}}\\[0.3em]
Anticipated questions and answers, organized by slide.

\section{Why $\alpha$ can't exceed 3}

\Q{Is this a proof?}
Yes, for $|\det M|=1$ matrices. The bound $\alpha \leq 1 + 2/(r-1)$
follows from $|\det M| = \lambda_1 \prod |\lambda_i| = 1$ and the
constraint $|\lambda_i| < 1$ for $i \geq 2$ (Pisot property).
The all-real case ($\alpha = 3$ for rank 2) is confirmed numerically
but not proven analytically.

\Q{Does this work for $|\det M| > 1$?}
The bound generalizes:
$\alpha \leq 1 + \frac{2}{r-1}(1 - \ln|\det M|/\ln\lambda_1)$.
For $|\det M| = 2$, rank 2: $\alpha \leq 3 - 2\ln 2/\ln\lambda_1 < 3$.

\Q{Could any 1D substitution tiling have $\alpha > 3$?}
Not with $|\det M| = 1$ (proven).
With $|\det M| > 1$, the bound is lower, not higher.
To get $\alpha > 3$ you would need $|\det M| < 1$ (fractional),
which is not possible for integer matrices.

\section{$\bar\Lambda$ ranking}

\Q{Is there a formula for $\bar\Lambda$ of substitution tilings?}
No general formula. For URL: $\bar\Lambda = (1+a^2)/6$ (exact).
For stealthy: $\bar\Lambda \approx 1/(\pi^2\chi)$.
For substitutions: $\bar\Lambda$ is determined by the full spatial
structure, not just the eigenvalues. We compute it numerically.

\Q{Why is Fibonacci ($\bar\Lambda = 0.201$) so close to the lattice ($1/6 = 0.167$)?}
Fibonacci has balanced tile lengths ($L/S = \varphi \approx 1.618$),
so the spacing variation is modest. The Fibonacci chain is one of
the most ``ordered'' quasicrystals.

\section{URL parameter sweep}

\Q{What happens at cloaking ($a = 1$)?}
The form factor $\tilde f(k) = \sin(\pi k)/(\pi k)$ vanishes at all
$k = 2\pi n$ ($n \neq 0$), eliminating all Bragg peaks.
The diffraction pattern becomes purely continuous.
$\bar\Lambda = 1/3$ exactly.

\Q{Is there a connection to the metallic means?}
Yes: $\bar\Lambda(\mu_n) \to 1/3$ as $n \to \infty$,
matching URL $a = 1$ exactly.
As $n$ grows, the metallic chain approaches a lattice with
rare quasi-random defects, which has the same disorder
as a uniformly displaced lattice.

\section{$g_2$-invariant patterns}

\Q{How is the lattice-gas different from URL $a=1$?}
It isn't, statistically. Lattice-gas: $j + U(0,1)$.
URL $a=1$: $j + U(-1/2, 1/2)$. Same distribution, shifted by $1/2$.

\Q{Why does Gaussian $\sigma = 0.2$ match Silver?}
Likely coincidence. The structures are completely different
(disordered Gaussian vs.\ deterministic quasicrystal),
but they produce the same large-scale fluctuation level.

\section{Metallic-mean convergence to $n = 100$}

\Q{Do any values overshoot $1/3$?}
Yes. $n = 50$: $\bar\Lambda = 0.334$; $n = 100$: $0.335$, both above $1/3 = 0.333$.
The trend is still upward. Error bars are $\sim 0.005$ but may be underestimated
at large $n$ (only 2 substitution iterations, $N = 300$k).
We cannot confirm convergence to exactly $1/3$ from this data.

\Q{So does it converge to $1/3$ or not?}
The $n = 1\ldots20$ data is clean and monotone from below.
The large-$n$ overshoots could be noise (poor statistics)
or could indicate a limit slightly above $1/3$.
Higher-precision runs at large $n$ (larger $N$) are needed to resolve this.

\Q{Is the convergence rate proven?}
No. The empirical fit $1/3 - \bar\Lambda \sim 0.22/n^{1.43}$
is from $n = 1\ldots20$ only and assumed convergence to $1/3$.

\section{Yuan-Torquato improved spreadability}

\Q{Why does this help for BT but not Fibonacci at high order?}
BT: $\alpha = 1.545$ (non-integer), spectral density has slower
oscillations $\to$ the asymptotic expansion converges at order 5--6.
Fibonacci: $\alpha = 3$ (integer), dense Bragg peaks cause the
expansion to diverge past order 1.

\Q{Is this better than variance-based $\bar\Lambda$?}
No. Variance converges at $N \sim 10^3$ with no fitting.
Yuan method is useful when you only have spreadability data
(e.g.\ NMR experiments) and need $\alpha$.

\Q{Did Yuan \& Torquato test on quasicrystals?}
No. Their Section VI explicitly notes that quasicrystalline media
with dense Bragg peaks are not covered by the regular expansion.

\section{$|\det M|{=}2$ chains}

\Q{Why 138 instead of 28?}
28 was from row sums $\leq 6$. Expanding to row sums $\leq 10$
found many more matrices satisfying the constraints.

\Q{Why does $\bar\Lambda$ vary so much within each $\alpha$ group?}
Matrices with the same eigenvalues (same $\alpha$) can have
very different substitution rules, producing different tile
arrangements. $\alpha$ captures only the leading spectral behavior;
$\bar\Lambda$ is sensitive to the full spatial structure.

\Q{Could we get $\alpha$ closer to 3?}
With $|\det M| = 2$: $\alpha = 3 - 2\ln 2/\ln\lambda_1$.
Need $\lambda_1 \to \infty$, meaning enormous row sums.
With $|\det M| = 3$: need $\lambda_1 > 9$.
At $\lambda_1 = 100$: $\alpha = 2.52$. Slow approach to 3.

\end{document}
